\documentclass{article}
\usepackage{amsmath,amssymb,amsthm}
\title{Ten Lectures on Integral Lattices}
\author{Lecture Notes}
\date{}
\newtheorem{theorem}{Theorem}
\newtheorem{definition}{Definition}
\begin{document}
\maketitle
\section{Lattices}
\begin{definition}
A lattice is a free $\mathbb{Z}$-module $L$ of finite rank with a symmetric bilinear form
$b \colon L \times L \to \mathbb{Q}$. It is integral if $b(L, L) \subseteq \mathbb{Z}$.
\end{definition}
\newpage
\section{Gram matrices}
For a basis $e_1, \dots, e_n$ of $L$ the Gram matrix is $G = (b(e_i, e_j))_{i,j}$, and
\[ \det L = \det G. \]
A change of basis by $P \in \mathrm{GL}_n(\mathbb{Z})$ replaces $G$ by $P^{T} G P$.
\newpage
\section{Dual lattices}
The dual lattice is $L^\vee = \{ x \in L \otimes \mathbb{Q} : b(x, L) \subseteq \mathbb{Z} \}$.
For integral $L$ we have $L \subseteq L^\vee$ and $[L^\vee : L] = |\det L|$.
\newpage
\section{Discriminant forms}
For an even lattice the discriminant group $A_L = L^\vee / L$ carries the quadratic form
\[ q_L(x + L) = \tfrac{1}{2} b(x, x) \bmod \mathbb{Z}. \]
\newpage
\section{Root lattices}
The root lattice $A_n$ is $\{ x \in \mathbb{Z}^{n+1} : \sum_i x_i = 0 \}$, with
$\det A_n = n + 1$. The lattice $D_n$ is $\{ x \in \mathbb{Z}^n : \sum_i x_i \equiv 0 \pmod 2 \}$.
\newpage
\section{The lattice $E_8$}
\begin{theorem}
$E_8 = D_8 \cup (D_8 + \tfrac{1}{2}(1, \dots, 1))$ is the unique even unimodular positive
definite lattice of rank $8$, and it has $240$ roots.
\end{theorem}
\newpage
\section{Theta series}
The theta series of a positive definite lattice is $\theta_L(q) = \sum_{x \in L} q^{b(x,x)/2}$.
For $E_8$,
\[ \theta_{E_8}(q) = 1 + 240 \sum_{n \geq 1} \sigma_3(n) q^n. \]
\newpage
\section{Genus}
Two lattices are in the same genus if $L \otimes \mathbb{Z}_p \cong M \otimes \mathbb{Z}_p$ for
every prime $p$ and $L \otimes \mathbb{R} \cong M \otimes \mathbb{R}$.
\newpage
\section{The mass formula}
The Smith--Minkowski--Siegel mass formula states
\[ \sum_{M \in \operatorname{gen}(L)} \frac{1}{|\mathrm{O}(M)|} = \operatorname{mass}(L). \]
\newpage
\section{Niemeier lattices}
There are exactly $24$ even unimodular positive definite lattices of rank $24$; the Leech
lattice is the only one without roots.
\end{document}
